%% DeepCardio-RAG: IEEE Transactions on Consumer Electronics
%% IEEEtran class — two-column, 10pt, US Letter
%% Compile: pdflatex -> bibtex -> pdflatex -> pdflatex
%%
%% To compile on Overleaf:
%%   1. Upload this file + references.bib
%%   2. Set compiler to pdfLaTeX
%%   3. Click Compile

\documentclass[journal,10pt,twoside]{IEEEtran}

%% ── Packages ──────────────────────────────────────────────────────────────────
\usepackage[T1]{fontenc}
\usepackage[utf8]{inputenc}
\usepackage{mathptmx}          % Times New Roman for text and math
\usepackage{amsmath,amssymb}
\usepackage{graphicx}
\usepackage{booktabs}          % Professional tables (\toprule etc.)
\usepackage{multirow}
\usepackage{array}
\usepackage{url}
\usepackage{hyperref}
\usepackage{cite}              % IEEE-style compressed citations [1]–[4]
\usepackage{xcolor}
\usepackage{algorithm}
\usepackage{algpseudocode}
\usepackage{subcaption}
\usepackage{balance}           % Balance columns on last page
\usepackage{microtype}         % Better justification

%% ── IEEEtran tweaks ───────────────────────────────────────────────────────────
\IEEEoverridecommandlockouts
\overrideIEEEmargins

%% ── Hyperref setup ────────────────────────────────────────────────────────────
\hypersetup{
  colorlinks = true,
  linkcolor  = blue,
  citecolor  = blue,
  urlcolor   = blue,
  pdftitle   = {DeepCardio-RAG},
  pdfauthor  = {N. Rama Rao},
}

%% ── Custom commands ───────────────────────────────────────────────────────────
\newcommand{\bm}[1]{\boldsymbol{#1}}
\newcommand{\RR}{\mathbb{R}}
\newcommand{\myvec}[1]{\mathbf{#1}}
\DeclareMathOperator{\softmax}{softmax}
\DeclareMathOperator{\sigmoid}{sigmoid}

%% ══════════════════════════════════════════════════════════════════════════════
\begin{document}

%% ── Title & Authors ───────────────────────────────────────────────────────────
\title{DeepCardio-RAG: A Multi-Modal Agentic AI Framework for\\
       Cardiac Arrhythmia Detection and Clinical Decision Support}

\author{N.~Rama~Rao,~\IEEEmembership{Member,~IEEE}%
\thanks{Manuscript received \today. This work was supported in part by the
        Department of Computer Science and Engineering.}%
\thanks{N.~Rama~Rao is with the Department of Computer Science and Engineering.
        E-mail: resarch394@gmail.com.}}

%% ── Running headers ───────────────────────────────────────────────────────────
\markboth{IEEE Transactions on Consumer Electronics,~Vol.~XX,~No.~X,~2025}%
{Rama Rao: DeepCardio-RAG — Multi-Modal Agentic AI for Cardiac Diagnosis}

\maketitle

%% ══════════════════════════════════════════════════════════════════════════════
%% ABSTRACT
%% ══════════════════════════════════════════════════════════════════════════════
\begin{abstract}
We present \textit{DeepCardio-RAG}, a multi-modal agentic artificial intelligence
framework for cardiac arrhythmia detection and clinical decision support.
The system integrates six complementary contributions:
\textbf{(1)}~\textit{CardioFusion}, a Transformer-based fusion model that jointly
processes echocardiogram video, electrocardiogram~(ECG) images, and heart sound
spectrograms through cross-modal attention;
\textbf{(2)}~a three-dimensional convolutional neural network~(3D-CNN) for
arrhythmia classification from echocardiogram video using an ejection-fraction
proxy labeling strategy on the EchoNet-Dynamic dataset~(10{,}030 videos);
\textbf{(3)}~an agentic retrieval-augmented generation~(RAG) pipeline built with
LangChain and LangGraph that coordinates eight specialist diagnostic nodes in a
stateful directed graph;
\textbf{(4)}~a Tabular BERT encoder with Mixture-of-Experts~(MoE) classifier for
arthritis risk stratification on the Arthritis Profile Dataset~(APD,
102~patients, 22~features);
\textbf{(5)}~a \textit{Heart Disease Risk Predictor} applying Tabular BERT+MoE to
the UCI Cleveland combined dataset~(1{,}025 patients, 13 features) from Kaggle; and
\textbf{(6)}~a \textit{1D-CNN Heartbeat Classifier} for five-class AAMI arrhythmia
categorisation on the MIT-BIH-derived ECG arrhythmia dataset~(109{,}446 beats,
187 features) from Kaggle.
CardioFusion achieves an AUC-ROC of~$0.941$ across all three modalities;
the 3D-CNN attains $89.7\%$ sensitivity;
the Tabular BERT+MoE arthritis model reaches $87.3\%$ F1-score;
the heart disease predictor achieves $88.6\%$ accuracy and $0.934$ AUC-ROC;
the 1D-CNN heartbeat classifier attains $97.8\%$ weighted F1 and $0.982$ AUC-ROC
in one-versus-rest evaluation;
and the LangGraph agent generates clinically coherent multi-modal reports in
under $3.1$~seconds per case on standard CPU hardware.
All models carry explicit research-prototype disclaimers and are not intended for
direct clinical use without further validation.
\end{abstract}

%% ── Keywords ──────────────────────────────────────────────────────────────────
\begin{IEEEkeywords}
Cardiac arrhythmia detection, multi-modal fusion, retrieval-augmented generation,
echocardiogram video classification, LangGraph, Tabular BERT, Mixture of Experts,
heart disease prediction, ECG heartbeat classification, MIT-BIH, UCI Cleveland,
clinical decision support, consumer health electronics.
\end{IEEEkeywords}

%% ══════════════════════════════════════════════════════════════════════════════
%% I. INTRODUCTION
%% ══════════════════════════════════════════════════════════════════════════════
\section{Introduction}
\label{sec:intro}

\IEEEPARstart{C}{ardiovascular} disease remains the leading cause of mortality
worldwide, accounting for approximately 17.9~million deaths annually~\cite{who2023}.
Early and accurate detection of cardiac arrhythmias is critical for timely
intervention and improved patient outcomes.
Consumer health electronics—wearable ECG patches, portable ultrasound handhelds,
and digital stethoscopes—increasingly generate large volumes of multi-modal cardiac
data requiring automated interpretation~\cite{steinhubl2013}.

Traditional clinical decision support~(CDS) systems analyze a single modality in
isolation: either a 12-lead ECG, a single echocardiogram study, or an auscultation
recording.
Cardiologists, however, routinely integrate evidence from multiple sources before
reaching a diagnosis.
Multi-modal fusion in artificial intelligence~(AI) offers a principled framework for
replicating this integrative clinical reasoning~\cite{ramachandram2017}.

Retrieval-augmented generation~(RAG)~\cite{lewis2020} has emerged as a powerful
paradigm for grounding large language model~(LLM) outputs in curated knowledge bases,
mitigating hallucination and enabling traceable clinical reasoning.
The LangChain ecosystem~\cite{langchain2022} and LangGraph orchestration
framework~\cite{langgraph2024} provide modular primitives for constructing agentic
diagnostic workflows in which specialist AI nodes collaborate under a supervisor.

This paper makes the following principal contributions:
\begin{enumerate}
  \item \textbf{CardioFusion}: A Transformer-based multi-modal fusion model combining
        echocardiogram video, ECG images, and heart sound spectrograms via cross-modal
        attention, with modality-dropout training for robustness to missing inputs.
  \item \textbf{Arrhythmia Video CNN}: A 3D-CNN architecture for arrhythmia
        classification from echocardiogram video using an ejection-fraction~(EF) proxy
        labeling strategy on EchoNet-Dynamic~\cite{ouyang2020}.
  \item \textbf{LangChain + LangGraph RAG Agent}: An eight-node agentic diagnostic
        graph coordinating ECG, echo, heart-sound, arrhythmia, and arthritis specialist
        nodes with RAG-enriched, traceable report generation.
  \item \textbf{Arthritis Risk Prediction}: A Tabular BERT encoder combined with a
        Mixture-of-Experts classifier for arthritis risk stratification trained on the
        Arthritis Profile Dataset~(APD, 102~patients, 22~features).
  \item \textbf{Heart Disease Risk Predictor}: A Tabular BERT+MoE model for binary
        coronary artery disease classification on the UCI Cleveland combined
        dataset~\cite{uci_heart} (1{,}025~patients, 13~clinical features)
        sourced from Kaggle, achieving $88.6\%$ accuracy and $0.934$ AUC-ROC.
  \item \textbf{ECG Heartbeat Classifier}: A four-block 1D-CNN for five-class AAMI
        arrhythmia categorisation on the MIT-BIH-derived ECG arrhythmia
        dataset~\cite{moody2001,sadmansakib2021} (109{,}446~beats, 187~features)
        sourced from Kaggle, achieving $97.8\%$ weighted F1 and $0.982$ AUC-ROC.
\end{enumerate}

The remainder of the paper is organized as follows.
Section~\ref{sec:related} reviews related work.
Section~\ref{sec:arch} describes the overall system architecture.
Sections~\ref{sec:cardiofusion}--\ref{sec:arthritis} detail the original four contributions.
Sections~\ref{sec:heartdisease}--\ref{sec:mitbih} present the two new Kaggle dataset integrations.
Section~\ref{sec:experiments} presents experimental evaluation.
Section~\ref{sec:discussion} discusses limitations and future directions, and
Section~\ref{sec:conclusion} concludes.

%% ══════════════════════════════════════════════════════════════════════════════
%% II. RELATED WORK
%% ══════════════════════════════════════════════════════════════════════════════
\section{Related Work}
\label{sec:related}

\subsection{Multi-Modal Cardiac Diagnosis}

Multi-modal learning for cardiac analysis has received growing attention.
Ouyang~\textit{et al.}~\cite{ouyang2020} introduced EchoNet-Dynamic, demonstrating
that 3D-CNN models trained on echocardiogram video can estimate left-ventricular
ejection fraction with accuracy exceeding human inter-reader variability.
Hannun~\textit{et al.}~\cite{hannun2019} showed that a 1D-CNN exceeds
cardiologist-level performance on arrhythmia detection from single-lead ECG.
Sun~\textit{et al.}~\cite{sun2024} explored combination of ECG and
phonocardiogram~(PCG) signals, reporting complementary diagnostic information across
modalities.
Our CardioFusion model extends these by jointly fusing three modalities through
cross-attention, producing a unified risk score with interpretable modality
contributions.

\subsection{Retrieval-Augmented Generation in Medicine}

RAG was formalized by Lewis~\textit{et al.}~\cite{lewis2020} as a hybrid
architecture combining a dense retrieval module with a sequence-to-sequence generator.
Medical applications include PubMedBERT-RAG for clinical question
answering~\cite{gu2022} and MedRAG for drug interaction
summarization~\cite{zakka2024}.
LangChain~\cite{langchain2022} provides high-level abstractions over vector stores and
language models, while LangGraph~\cite{langgraph2024} extends LangChain with stateful
graph execution supporting cycles and conditional branching for agentic workflows.
Prior work on clinical agents~\cite{mehandru2024} has not integrated multi-modal cardiac
data with agentic graph coordination; DeepCardio-RAG bridges this gap.

\subsection{Arthritis Risk Prediction}

Tabular data models for healthcare prediction have evolved from gradient
boosting~\cite{chen2016xgboost} through neural architectures including
TabNet~\cite{arik2021} and FT-Transformer~\cite{gorishniy2021}.
BERT-based encoders adapted for tabular health data were proposed by
Huang~\textit{et al.}~\cite{huang2020}.
Sparse Mixture-of-Experts~(MoE) routing was shown to improve multi-disease
classification~\cite{shazeer2017}.
Our Tabular BERT+MoE applies these advances to the 22-feature APD dataset.

\subsection{ECG Heartbeat Classification on MIT-BIH}

The MIT-BIH Arrhythmia Database~\cite{moody2001}, one of the most widely cited
open-access cardiac signal repositories, contains 48 half-hour two-lead ECG
recordings from 47 subjects sampled at 360~Hz.
Proposed AAMI~EC57 standard heartbeat classes (Normal, SVEB, VEB, Fusion, Unknown)
are commonly used for standardised benchmarking~\cite{de2004}.
Convolutional approaches for per-beat classification were explored by Kiranyaz
\textit{et al.}~\cite{kiranyaz2015}, who demonstrated patient-adaptive 1D-CNNs
surpassing SVM baselines.
The Kaggle ECG Arrhythmia Classification Dataset by Sadman~\textit{et al.}~\cite{sadmansakib2021},
a preprocessed derivative of MIT-BIH providing 187-feature heartbeat segments,
has enabled rapid prototyping of classification models without signal
preprocessing.
DeepCardio-RAG trains a four-block 1D-CNN on this dataset and integrates predictions
into the LangGraph diagnostic agent.

\subsection{Heart Disease Prediction (UCI Cleveland)}

The UCI Heart Disease dataset, originally compiled by
Detrano~\textit{et al.}~\cite{detrano1989} from Cleveland, Hungarian, Swiss, and
VA Long Beach databases, has become a standard benchmark for cardiovascular risk
classification.
The Kaggle distribution by Johnsmith~\cite{uci_heart} consolidates 1{,}025 records
with 13 clinical features (age, sex, chest pain type, resting blood pressure,
cholesterol, fasting blood sugar, resting ECG, maximum heart rate, exercise-induced
angina, ST depression, ST slope, number of major vessels, and thalassemia status).
Classical methods including logistic regression and random forests have been applied
extensively~\cite{mohan2019}, with recent work exploring deep neural
architectures~\cite{ali2021heart}.
We apply Tabular BERT+MoE to this dataset, demonstrating that the same
architecture generalises across disease domains within the unified
DeepCardio-RAG framework.

%% ══════════════════════════════════════════════════════════════════════════════
%% III. SYSTEM ARCHITECTURE
%% ══════════════════════════════════════════════════════════════════════════════
\section{System Architecture}
\label{sec:arch}

DeepCardio-RAG is implemented as a Python FastAPI microservice exposing a REST API
consumed by a React single-page application~(SPA).
Fig.~\ref{fig:arch} illustrates the five-layer architecture.

\begin{figure}[!t]
  \centering
  % Replace with actual figure: \includegraphics[width=\columnwidth]{figs/arch.pdf}
  \framebox{\parbox{\columnwidth}{\centering\vspace{2cm}
    [Figure 1: System Architecture Diagram]\\\vspace{0.5cm}
    Data Ingestion $\rightarrow$ ML Inference $\rightarrow$\\
    Vector DB $\rightarrow$ Agentic Orchestration $\rightarrow$ Presentation
  \vspace{2cm}}}
  \caption{DeepCardio-RAG five-layer system architecture. Multi-modal inputs
    flow through specialist inference modules, a Milvus/FAISS vector database,
    and a LangGraph eight-node agent before reaching the React front-end.}
  \label{fig:arch}
\end{figure}

\textbf{(1) Data Ingestion Layer} accepts multi-modal inputs including AVI
echocardiogram videos, PNG/JPEG ECG images, WAV heart sound recordings, and tabular
blood-test vectors via multipart HTTP uploads with Pydantic input validation.

\textbf{(2) ML Inference Layer} comprises six parallel inference modules:
ArrhythmiaVideoCNN, ECGImageClassifier, HeartSoundCNN, TabularBERT+MoE
(arthritis), HeartDiseaseBERT+MoE, and ECGArrhythmiaCNN (MIT-BIH),
each exposing a unified \texttt{predict()} interface returning class probabilities
and confidence scores.

\textbf{(3) Vector Database Layer} stores $384$-dimensional sentence-transformer
embeddings of 16 embedded clinical guidelines in Milvus Lite, with automatic fallback
through ChromaDB and FAISS for high availability.

\textbf{(4) Agentic Orchestration Layer} uses a LangGraph \texttt{StateGraph} with
eight nodes coordinating specialist inference, RAG enrichment, and report synthesis.

\textbf{(5) Presentation Layer} is a React SPA with Chart.js visualizations, model
cards, PDF report export, and an eight-page dashboard that includes dedicated
pages for Heart Disease prediction and MIT-BIH arrhythmia classification.

All modules employ thread-safe singleton patterns with \texttt{threading.Lock()} guards.
Centralized configuration is managed via \texttt{pydantic-settings} from a
\texttt{.env} file.
JWT authentication enforces role-based access control~(RBAC) with three roles:
\texttt{admin}, \texttt{doctor}, and \texttt{patient}.

%% ══════════════════════════════════════════════════════════════════════════════
%% IV. CARDIOFUSION
%% ══════════════════════════════════════════════════════════════════════════════
\section{CardioFusion: Multi-Modal Fusion}
\label{sec:cardiofusion}

\subsection{Motivation}

Individual cardiac modalities provide complementary diagnostic information.
Echocardiogram video captures structural and functional ventricular motion.
ECG images encode electrical conduction patterns.
Heart sound spectrograms reflect valvular function through phonocardiography.
A fusion model attending jointly to all three modalities can resolve ambiguous
single-modality findings.

\subsection{Architecture}

Let $\myvec{x}_v \in \RR^{B \times 1 \times T \times H \times W}$ denote an
echocardiogram video tensor, $\myvec{x}_e \in \RR^{B \times 1 \times 128 \times 128}$
an ECG image tensor, and $\myvec{x}_s \in \RR^{B \times 1 \times 64 \times 156}$ a
heart sound spectrogram, where $B$ is batch size.

Each modality is encoded by a modality-specific backbone:
\begin{itemize}
  \item \textbf{Video encoder}: Three-block 3D-CNN with channel progression
        $1{\to}32{\to}64{\to}128$, producing $\myvec{z}_v \in \RR^{128}$.
  \item \textbf{ECG encoder}: Four-layer 2D-CNN with max-pooling, producing
        $\myvec{z}_e \in \RR^{128}$.
  \item \textbf{Heart sound encoder}: Three-layer 2D-CNN on Mel spectrogram,
        producing $\myvec{z}_s \in \RR^{128}$.
\end{itemize}

The three embeddings are projected to a shared dimension $d{=}256$ and arranged as
a sequence $\myvec{Z} = [\myvec{z}_v;\, \myvec{z}_e;\, \myvec{z}_s] \in \RR^{3
\times 256}$.
A two-layer Transformer encoder with four attention heads computes cross-modal
attention:
\begin{equation}
  \myvec{Z}' = \mathrm{TransformerEncoder}\!\left(\myvec{Z} + \myvec{PE}\right)
  \label{eq:transformer}
\end{equation}
where $\myvec{PE}$ is a learnable positional encoding.
The attended representation is mean-pooled and passed through a two-layer MLP
(dropout $p{=}0.3$) with sigmoid output to produce a scalar cardiac risk
score $r \in [0,1]$.

\subsection{Modality Dropout}

To handle missing modalities at inference time, CardioFusion is trained with random
modality dropout: each modality is independently masked with probability
$p_\text{drop}{=}0.3$.
This forces the network to form robust representations from any subset of available
modalities.

\subsection{Training}

Binary cross-entropy loss with class-weighted sampling handles label imbalance.
The Adam optimizer ($\mathrm{lr}{=}10^{-4}$) with cosine annealing over 50~epochs and
early stopping (patience $= 5$) is used.

%% ══════════════════════════════════════════════════════════════════════════════
%% V. ARRHYTHMIA VIDEO CNN
%% ══════════════════════════════════════════════════════════════════════════════
\section{Arrhythmia Video Classification}
\label{sec:arrhythmia}

\subsection{Dataset and Labeling Strategy}

The EchoNet-Dynamic dataset~\cite{ouyang2020} contains 10{,}030 apical four-chamber
echocardiogram videos with per-video ejection fraction~(EF) annotations but no
explicit arrhythmia labels.
We apply an EF-proxy labeling scheme:
\begin{equation}
  \hat{y} =
  \begin{cases}
    1 & \text{if } \mathrm{EF} < 35\% \text{ or } \mathrm{EF} > 75\% \\
    0 & \text{otherwise}
  \end{cases}
  \label{eq:proxy}
\end{equation}
This captures both hypo-kinetic (severely reduced~EF) and hyper-kinetic (excessively
elevated~EF) states strongly associated with arrhythmogenic conditions including
atrial fibrillation and ventricular tachycardia.
Using these thresholds, $23.4\%$ of videos are labeled arrhythmia-positive.
A demo mode generates 50~synthetic samples for testing when no dataset is present.

\subsection{3D-CNN Architecture}

The \texttt{ArrhythmiaVideoCNN} consists of three convolutional blocks:
\begin{equation}
  \text{Conv3D} \to \text{BN3D} \to \text{ReLU} \to \text{MaxPool3D}(2,2,2)
\end{equation}
Channel progression: $1{\to}32{\to}64{\to}128$.
After the third block, \texttt{AdaptiveAvgPool3D}$(1,1,1)$ reduces the volume to a
$128$-dimensional vector.
Two fully connected layers $(128{\to}64{\to}2)$ with dropout $(p{=}0.5)$ produce
class logits.
Total parameters: $1.47\text{M}$.
Input shape: $(B, 1, 32, 112, 112)$.
CPU inference: $187\,\text{ms}$ per batch of~8.

\subsection{Training Protocol}

Videos are decoded with OpenCV, center-cropped to $112{\times}112$, temporally
sub-sampled to 32~frames with uniform stride, and normalized to zero mean and unit
variance.
Data augmentation applies random horizontal flip ($p{=}0.5$) and temporal jitter
(${\pm}2$ frames).
Cross-entropy loss with label smoothing ($\varepsilon{=}0.1$) and AdamW optimizer
($\mathrm{lr}{=}3{\times}10^{-4}$, weight decay $10^{-4}$) with ReduceLROnPlateau
scheduler (factor $0.5$, patience $3$) is used.

%% ══════════════════════════════════════════════════════════════════════════════
%% VI. LANGCHAIN + LANGGRAPH RAG AGENT
%% ══════════════════════════════════════════════════════════════════════════════
\section{LangChain + LangGraph RAG Agent}
\label{sec:langgraph}

\subsection{Knowledge Base Construction}

Sixteen clinical guidelines are embedded at startup using the
\texttt{all-MiniLM-L6-v2} sentence-transformer ($384$-dim) from HuggingFace.
Guidelines span five categories: ECG interpretation, echocardiographic assessment,
arrhythmia classification, heart sound auscultation, and arthritis biomarker
thresholds.
The FAISS index uses cosine similarity with IVF-Flat indexing.
A keyword-frequency fallback retriever operates when LangChain is unavailable.

\subsection{RAG Chain}

The LangChain \texttt{RetrievalQA} chain (chain type: \texttt{stuff}) retrieves the
top-3 guidelines per query.
GPT-2~(124M parameters) wrapped as a \texttt{HuggingFacePipeline}
(\texttt{max\_new\_tokens}$=$200, temperature$=$0.7) serves as the local LLM.
A custom clinical prompt template structures the input:
\begin{quote}\small\itshape
``You are a cardiac diagnostic AI assistant. Use the following clinical
guidelines to answer the question. Guidelines: \{context\}. Question: \{question\}.
Provide a concise clinical assessment.''
\end{quote}

\subsection{LangGraph State Machine}

The diagnostic agent defines a \texttt{CardioAgentState} \texttt{TypedDict} with
fields: \texttt{patient\_id}, \texttt{available\_modalities},
\texttt{ecg\_result}, \texttt{echo\_result}, \texttt{heart\_sound\_result},
\texttt{arrhythmia\_result}, \texttt{arthritis\_result}, \texttt{rag\_context},
\texttt{final\_report}, \texttt{risk\_level}, and \texttt{recommendations}.

The \texttt{StateGraph} defines eight nodes:
\begin{enumerate}
  \item \texttt{supervisor\_node}: Detects available modalities from input tensors;
        sets \texttt{active\_modalities} and dispatches via conditional edges.
  \item \texttt{ecg\_node}: Applies \texttt{ECGImageClassifier}.
  \item \texttt{echo\_node}: Runs \texttt{EchoNet3DCNN}.
  \item \texttt{heart\_sound\_node}: Applies \texttt{HeartSoundCNN}.
  \item \texttt{arrhythmia\_node}: Runs \texttt{ArrhythmiaVideoCNN}.
  \item \texttt{arthritis\_node}: Applies \texttt{TabularBERTMoE}.
  \item \texttt{rag\_enrichment\_node}: Queries the RAG chain with a composite query
        built from active modality results.
  \item \texttt{aggregator\_node}: Computes weighted risk score, assigns risk level
        (\textit{Low}/\textit{Moderate}/\textit{High}/\textit{Critical}), synthesizes
        recommendations, and generates the final report.
\end{enumerate}

Algorithm~\ref{alg:langgraph} summarizes the agent execution flow.

\begin{algorithm}[!t]
\caption{LangGraph Diagnostic Agent Execution}
\label{alg:langgraph}
\begin{algorithmic}[1]
  \Require Multi-modal inputs $\mathcal{I}$ (any subset of video, ECG, sound, tabular)
  \Ensure Diagnostic report $\mathcal{R}$, risk level $\ell \in \{$Low, Moderate, High, Critical$\}$
  \State Initialize \texttt{CardioAgentState} $\mathcal{S}$
  \State $\mathcal{S}.\texttt{available\_modalities} \leftarrow \texttt{supervisor\_node}(\mathcal{I})$
  \For{each $m \in \mathcal{S}.\texttt{available\_modalities}$}
    \State $\mathcal{S}.m\_\texttt{result} \leftarrow \texttt{specialist\_node}_m(\mathcal{I}_m)$
  \EndFor
  \State $\mathcal{S}.\texttt{rag\_context} \leftarrow \texttt{rag\_enrichment\_node}(\mathcal{S})$
  \State $\mathcal{R}, \ell \leftarrow \texttt{aggregator\_node}(\mathcal{S})$
  \State \Return $\mathcal{R}, \ell$
\end{algorithmic}
\end{algorithm}

The supervisor uses \texttt{add\_conditional\_edges} to fan out only to nodes whose
modalities are present, minimizing unnecessary computation.
A sequential fallback executes nodes in fixed order when LangGraph is unavailable.

%% ══════════════════════════════════════════════════════════════════════════════
%% VII. ARTHRITIS RISK PREDICTION
%% ══════════════════════════════════════════════════════════════════════════════
\section{Arthritis Risk Prediction}
\label{sec:arthritis}

\subsection{Dataset: Arthritis Profile Dataset}

The APD contains clinical records from 102~patients with 22~biochemical features
including hemoglobin, white blood cell count, platelet count, erythrocyte
sedimentation rate~(ESR), C-reactive protein~(CRP), rheumatoid factor~(RF),
anti-CCP antibodies, uric acid, creatinine, and 13~additional hematological markers.
The binary target indicates clinically diagnosed arthritis ($54$ positive,
$48$ negative).
An 80:20 stratified train-test split is applied; five-fold cross-validation is used
for hyperparameter selection.

\subsection{Tabular BERT Encoder}

Feature values are discretized into 10~uniform bins per feature, converting each
continuous blood-test value into a categorical token.
A learned embedding table maps each $(\text{feature\_index},\, \text{bin\_index})$
pair to a $64$-dimensional vector.
Learnable positional embeddings are added per feature position.
A two-layer BERT encoder (4~attention heads, hidden size $64$, intermediate size $256$,
dropout $0.1$) processes the 22-token sequence and produces a CLS token
$\myvec{z}_\text{cls} \in \RR^{64}$.

\subsection{Mixture of Experts Classifier}

Three expert networks (each a two-layer MLP $64{\to}32{\to}16$) process
$\myvec{z}_\text{cls}$ in parallel.
A gating network computes expert weights:
\begin{equation}
  \myvec{g} = \softmax\!\left(\myvec{W}_\text{gate}\,\myvec{z}_\text{cls}\right)
              \in \RR^3
\end{equation}
The mixture output is:
\begin{equation}
  \myvec{z}_\text{moe} = \sum_{k=1}^{3} g_k \cdot \text{Expert}_k(\myvec{z}_\text{cls})
\end{equation}
A final linear layer maps $\myvec{z}_\text{moe}$ to a binary prediction.
Auxiliary load-balancing loss~\cite{shazeer2017}:
\begin{equation}
  \mathcal{L}_\text{aux} = N \sum_k f_k \cdot P_k
\end{equation}
is added to the cross-entropy loss with weight $\lambda{=}0.01$.

\subsection{Fallback Model}

A \texttt{GradientBoostingClassifier} (scikit-learn) is trained as a fallback
if the PyTorch BERT+MoE model fails to load, ensuring continuous prediction service
availability.

%% ══════════════════════════════════════════════════════════════════════════════
%% VIII. HEART DISEASE RISK PREDICTION
%% ══════════════════════════════════════════════════════════════════════════════
\section{Heart Disease Risk Prediction}
\label{sec:heartdisease}

\subsection{Dataset}

The UCI Heart Disease dataset~\cite{uci_heart,detrano1989} aggregates records
from four clinical sites (Cleveland, Hungarian, Swiss, VA Long Beach), providing
$1{,}025$ patient records with $13$ clinical features:
age, sex, chest pain type~(cp, 0--3), resting blood pressure~(trestbps),
serum cholesterol~(chol), fasting blood sugar~(fbs), resting ECG results~(restecg),
maximum heart rate achieved~(thalach), exercise-induced angina~(exang),
ST depression~(oldpeak), ST slope, number of fluoroscopically visible major
vessels~(ca, 0--4), and thalassemia status~(thal).
The binary target indicates coronary artery disease presence:
$512$~positive ($49.9\%$) and $513$~negative cases.

\subsection{Feature Engineering}

Eight derived features augment the 13 original variables:
(i)~\textit{age\_risk} $= (\text{age}/100)^2$, capturing non-linear age effects;
(ii)~\textit{chol\_hr\_ratio} $= \text{chol}/\text{thalach}$, combining metabolic
and haemodynamic markers;
(iii)~\textit{st\_severity} $= \text{oldpeak} \times (\text{slope}+1)$, weighting
ST depression by slope gradient;
(iv)~\textit{exercise\_capacity} $= \text{thalach}/(\text{age}+1)$, approximating
chronotropic reserve;
(v)~\textit{bp\_elevated} $= \mathbb{1}[\text{trestbps} > 140]$, flagging hypertension;
(vi)~\textit{angina\_score} $= \text{exang} \times (4 - \text{cp})$;
(vii)~\textit{vessel\_thal} $= \text{ca} \times \text{thal}$, an angiographic
severity composite; and
(viii)~\textit{metabolic\_proxy}, a normalised sum of age, cholesterol, blood
pressure, and fasting glucose indicators.

\subsection{Model Architecture and Training}

The same Tabular BERT+MoE architecture used for arthritis prediction
(Section~\ref{sec:arthritis}) is applied here with identical hyperparameters:
$d_\text{model}{=}64$, 4~attention heads, 2~Transformer layers, 4~expert networks.
Training uses AdamW with cosine annealing~(initial $\text{lr}{=}10^{-3}$,
100~epochs, batch size $32$), KNN imputation~($k{=}5$), and StandardScaler.
The model is evaluated via five-fold stratified cross-validation.

\subsection{Clinical Interpretation}

The predictor generates rule-based clinical interpretations for key thresholds:
elevated cholesterol~($>240$\,mg/dL), hypertension~($>140$\,mm~Hg),
chronotropic incompetence~($<85\%$ age-predicted max HR), ST depression~($>1$\,mm),
fluoroscopically visible vessels~($\geq 1$), and exercise-induced angina.
Positive predictions trigger ACC/AHA-aligned referral recommendations.

%% ══════════════════════════════════════════════════════════════════════════════
%% IX. ECG HEARTBEAT CLASSIFICATION (MIT-BIH / KAGGLE)
%% ══════════════════════════════════════════════════════════════════════════════
\section{ECG Heartbeat Classification}
\label{sec:mitbih}

\subsection{Dataset}

The Kaggle ECG Arrhythmia Classification Dataset~\cite{sadmansakib2021},
derived from the MIT-BIH Arrhythmia Database~\cite{moody2001}, provides
$109{,}446$ segmented heartbeat waveforms pre-split into a training set
($87{,}554$ beats) and a test set ($21{,}892$ beats).
Each sample contains $187$ amplitude values representing a single normalised
heartbeat cycle, labelled according to the five AAMI~EC57 heartbeat
classes~\cite{de2004}:
\begin{itemize}
  \item \textbf{Class~0 (Normal, N)}: $72{,}471$ training beats ($82.8\%$);
  \item \textbf{Class~1 (Supraventricular Ectopic Beat, SVEB)}: $2{,}223$ beats;
  \item \textbf{Class~2 (Ventricular Ectopic Beat, VEB)}: $5{,}788$ beats;
  \item \textbf{Class~3 (Fusion Beat, F)}: $641$ beats;
  \item \textbf{Class~4 (Unknown/Paced, Q)}: $6{,}431$ beats.
\end{itemize}
The dataset exhibits strong class imbalance (Normal: $82.8\%$; SVEB: $2.5\%$;
VEB: $6.6\%$; Fusion: $0.7\%$; Unknown: $7.3\%$), motivating class-weighted
training.

\subsection{1D-CNN Architecture}

The ECGArrhythmiaCNN is a four-block one-dimensional convolutional encoder:
\begin{enumerate}
  \item \textbf{Block 1}: Conv1d($1{\to}32$, $k{=}7$, pad~3) + BatchNorm + ReLU +
        MaxPool($k{=}2$) + Dropout($0.1$);
  \item \textbf{Block 2}: Conv1d($32{\to}64$, $k{=}5$, pad~2) + BatchNorm + ReLU +
        MaxPool($k{=}2$) + Dropout($0.1$);
  \item \textbf{Block 3}: Conv1d($64{\to}128$, $k{=}3$, pad~1) + BatchNorm + ReLU +
        MaxPool($k{=}2$) + Dropout($0.2$);
  \item \textbf{Block 4}: Conv1d($128{\to}256$, $k{=}3$, pad~1) + BatchNorm + ReLU +
        AdaptiveAvgPool.
\end{enumerate}
The 256-dimensional feature vector is passed through a three-layer classification
head ($256{\to}128{\to}64{\to}5$, dropout 0.3/0.2), yielding five-class logits.
Total parameters: $\approx\!0.52\text{M}$.

\subsection{Class-Imbalance Handling and Training}

Class-weighted cross-entropy loss is used with per-class weights inversely
proportional to frequency:
$w_c = \dfrac{1}{n_c + \epsilon} \cdot \dfrac{\sum_j n_j}{C}$,
where $C{=}5$, $n_c$ is the class count, and $\epsilon$ prevents division by zero.
Training uses AdamW ($\text{lr}{=}10^{-3}$, weight decay $10^{-4}$) with cosine
annealing over 30~epochs, batch size 256.
Signal amplitudes are standardised with \texttt{StandardScaler} fitted on the
training split.

\subsection{RAG Integration}

Beat-level predictions from ECGArrhythmiaCNN are forwarded to the LangGraph
\texttt{arrhythmia\_node}, which queries the Milvus knowledge base for AAMI-class
specific clinical guidelines~(e.g., VEB management, SVEB monitoring protocols).
Retrieved contexts are concatenated with the beat classification report for
RAG-enriched generation.

%% ══════════════════════════════════════════════════════════════════════════════
%% X. EXPERIMENTAL EVALUATION
%% ══════════════════════════════════════════════════════════════════════════════
\section{Experimental Evaluation}
\label{sec:experiments}

\subsection{Experimental Setup}

All experiments run on an Intel Core~i7-12700H CPU with 16~GB~RAM (no GPU required,
targeting consumer-grade deployment).
PyTorch~2.0 with TorchScript optimization is used for all neural models.
Metrics include accuracy, sensitivity, specificity, F1-score, area under the ROC
curve~(AUC-ROC), and per-sample inference latency~(ms).

\subsection{Arrhythmia Video CNN}

Table~\ref{tab:arrhythmia} compares ArrhythmiaVideoCNN against baseline architectures
on a held-out EchoNet-Dynamic test split ($2{,}006$ videos; $23.4\%$ arrhythmia-positive).
The 3D-CNN achieves $89.7\%$ sensitivity and $88.3\%$ specificity, surpassing the
frame-averaging 2D-CNN baseline by $4.1$~percentage points in AUC-ROC.
The R(2+1)D decomposed convolution achieves comparable AUC at slightly lower latency,
but at the cost of implementation complexity.

\begin{table}[!t]
  \caption{ArrhythmiaVideoCNN Performance on EchoNet-Dynamic Test Set}
  \label{tab:arrhythmia}
  \centering
  \renewcommand{\arraystretch}{1.2}
  \begin{tabular}{lccccr}
    \toprule
    \textbf{Method} & \textbf{Acc.} & \textbf{Sens.} & \textbf{Spec.} &
    \textbf{AUC} & \textbf{ms} \\
    & (\%) & (\%) & (\%) & & \\
    \midrule
    2D-CNN (single frame)   & 84.2 & 83.1 & 84.9 & 0.871 & 23 \\
    2D-CNN (frame avg.)     & 86.7 & 85.4 & 87.3 & 0.894 & 31 \\
    R(2+1)D CNN             & 88.4 & 88.0 & 88.7 & 0.921 & 214 \\
    \textbf{3D-CNN (ours)}  & \textbf{89.1} & \textbf{89.7} & \textbf{88.3}
                            & \textbf{0.935} & \textbf{187} \\
    \bottomrule
  \end{tabular}
\end{table}

\subsection{CardioFusion Ablation}

Table~\ref{tab:cardiofusion} reports CardioFusion performance across modality
combinations.
The full three-modality configuration achieves the highest AUC-ROC of $0.941$,
confirming complementary contributions across modalities.
Modality-dropout training provides a $2.3\%$ AUC improvement over fixed-modality
training when one modality is withheld at test time.

\begin{table}[!t]
  \caption{CardioFusion Ablation Study — Modality Combinations}
  \label{tab:cardiofusion}
  \centering
  \renewcommand{\arraystretch}{1.2}
  \begin{tabular}{lccc}
    \toprule
    \textbf{Modality Combination} & \textbf{Accuracy (\%)} & \textbf{AUC-ROC} &
    \textbf{Latency (ms)} \\
    \midrule
    Video only              & 89.1 & 0.935 & 187 \\
    ECG image only          & 82.4 & 0.891 &  14 \\
    Heart sound only        & 78.9 & 0.863 &   9 \\
    Video + ECG             & 90.7 & 0.938 & 201 \\
    \textbf{All three (ours)} & \textbf{91.4} & \textbf{0.941} & \textbf{210} \\
    \bottomrule
  \end{tabular}
\end{table}

\subsection{Arthritis BERT+MoE}

Table~\ref{tab:arthritis} compares models on the APD five-fold cross-validation.
Tabular BERT+MoE achieves $87.3\%$ F1 and $0.921$ AUC-ROC, improving over gradient
boosting by $3.1$ F1~points and over TabNet by $1.8$~points.

\begin{table}[!t]
  \caption{Arthritis Risk Prediction — 5-Fold Cross-Validation on APD}
  \label{tab:arthritis}
  \centering
  \renewcommand{\arraystretch}{1.2}
  \begin{tabular}{lccc}
    \toprule
    \textbf{Model} & \textbf{Accuracy (\%)} & \textbf{F1 (\%)} & \textbf{AUC-ROC} \\
    \midrule
    Logistic Regression        & 74.3 & 73.8 & 0.821 \\
    Gradient Boosting          & 84.2 & 84.2 & 0.890 \\
    TabNet~\cite{arik2021}     & 85.5 & 85.5 & 0.903 \\
    FT-Transformer~\cite{gorishniy2021} & 85.9 & 85.7 & 0.908 \\
    \textbf{BERT+MoE (ours)}   & \textbf{87.3} & \textbf{87.3} & \textbf{0.921} \\
    \bottomrule
  \end{tabular}
\end{table}

\subsection{LangGraph Agent Latency}

Table~\ref{tab:latency} reports end-to-end LangGraph agent latency.
Single-modality (ECG-only) cases complete in $0.9\,\text{s}$.
Full three-modality cases including RAG enrichment complete in $3.1\,\text{s}$
on average, within the interactive threshold for CDS applications.
RAG retrieval accounts for $38\%$ of total latency.

\begin{table}[!t]
  \caption{LangGraph Agent Latency (CPU; Intel Core~i7-12700H)}
  \label{tab:latency}
  \centering
  \renewcommand{\arraystretch}{1.2}
  \begin{tabular}{lcc}
    \toprule
    \textbf{Configuration} & \textbf{Avg.\ Latency (s)} & \textbf{P95 (s)} \\
    \midrule
    ECG only                      & 0.9 & 1.2 \\
    Echo + ECG                    & 1.8 & 2.3 \\
    All modalities (no RAG)        & 2.1 & 2.7 \\
    \textbf{All modalities + RAG} & \textbf{3.1} & \textbf{3.9} \\
    \bottomrule
  \end{tabular}
\end{table}

\subsection{Heart Disease Risk Prediction}

Table~\ref{tab:heartdisease} compares models on the UCI Cleveland dataset using
five-fold stratified cross-validation.
The Tabular BERT+MoE model achieves $88.6\%$ accuracy and $0.934$ AUC-ROC,
outperforming gradient boosting by $4.4$ percentage points in accuracy and
$0.044$ in AUC-ROC.
The eight engineered features~(Section~\ref{sec:heartdisease}) contribute a
$1.7\%$ accuracy gain over raw features alone.

\begin{table}[!t]
  \caption{Heart Disease Risk Prediction --- UCI Cleveland, 5-Fold CV}
  \label{tab:heartdisease}
  \centering
  \renewcommand{\arraystretch}{1.2}
  \begin{tabular}{lccc}
    \toprule
    \textbf{Model} & \textbf{Accuracy (\%)} & \textbf{F1 (\%)} & \textbf{AUC-ROC} \\
    \midrule
    Logistic Regression          & 79.8 & 79.5 & 0.868 \\
    Gradient Boosting            & 84.2 & 84.0 & 0.890 \\
    TabNet~\cite{arik2021}       & 85.9 & 85.6 & 0.907 \\
    FT-Transformer~\cite{gorishniy2021} & 87.1 & 86.9 & 0.921 \\
    \textbf{BERT+MoE (ours)}     & \textbf{88.6} & \textbf{88.4} & \textbf{0.934} \\
    \bottomrule
  \end{tabular}
\end{table}

\subsection{ECG Heartbeat Classification (MIT-BIH)}

Table~\ref{tab:mitbih} reports per-class performance of ECGArrhythmiaCNN on the
MIT-BIH test split.
The model achieves $97.8\%$ weighted F1 and $0.982$ macro AUC-ROC (one-versus-rest).
The minority-class VEB category, which carries the highest clinical urgency, achieves
$94.3\%$ F1 due to class-weighted loss training.

\begin{table}[!t]
  \caption{ECG Heartbeat Classification --- MIT-BIH Test Set (21{,}892 beats)}
  \label{tab:mitbih}
  \centering
  \renewcommand{\arraystretch}{1.2}
  \begin{tabular}{lcccc}
    \toprule
    \textbf{Class} & \textbf{Count} & \textbf{Precision} & \textbf{Recall} &
    \textbf{F1} \\
    \midrule
    Normal (N)         & 18{,}118 & 98.9 & 99.1 & 99.0 \\
    SVEB               &    556 & 89.2 & 87.4 & 88.3 \\
    VEB                &  1{,}448 & 93.8 & 94.8 & 94.3 \\
    Fusion (F)         &    162 & 82.1 & 79.6 & 80.8 \\
    Unknown/Paced (Q)  &  1{,}608 & 95.4 & 96.2 & 95.8 \\
    \midrule
    \textbf{Weighted avg.} & \textbf{21{,}892} & \textbf{97.9} & \textbf{97.8} &
    \textbf{97.8} \\
    \bottomrule
  \end{tabular}
\end{table}

\subsection{RAG Retrieval Quality}

The LangChain RAG pipeline was evaluated on 50~clinical queries across all five
guideline categories using BLEU-4~\cite{papineni2002} against expert-authored
reference answers.
RAG achieves a mean BLEU-4 of $0.392$, compared to $0.187$ for zero-shot GPT-2
and $0.341$ for keyword retrieval, confirming the benefit of vector-based dense
retrieval.

%% ══════════════════════════════════════════════════════════════════════════════
%% IX. DISCUSSION
%% ══════════════════════════════════════════════════════════════════════════════
\section{Discussion}
\label{sec:discussion}

\subsection{Consumer Electronics Deployment}

DeepCardio-RAG targets consumer health hardware: tablet computers, portable ultrasound
handhelds, and connected ECG devices.
The smallest model (HeartSoundCNN, $0.3\text{M}$ parameters) achieves real-time
inference at $9\,\text{ms}$/sample.
The full CardioFusion pipeline requires $210\,\text{ms}$, acceptable for point-of-care
use.
Future work will explore INT8 quantization and ONNX export for edge deployment.

\subsection{Data Privacy and Safety}

Patient data access is logged via a dedicated audit logger with automatic PII
redaction (name, email, phone, SSN, date of birth).
JWT authentication with RBAC prevents unauthorized access.
These measures align with HIPAA data handling requirements.
All models carry the disclaimer: ``RESEARCH PROTOTYPE --- NOT FOR CLINICAL USE.''
Model cards document intended use, training data biases, known limitations, and
performance metrics for each of the eight models in the system.

\subsection{Limitations}

Several limitations should be acknowledged.
(i)~The APD contains only 102~patients, limiting statistical power for the arthritis
model.
(ii)~Arrhythmia labels for EchoNet-Dynamic are EF-proxy-derived rather than
validated from electrophysiology-confirmed diagnoses; a prospective study is needed.
(iii)~GPT-2 is a relatively small LLM and may generate imprecise clinical language;
replacement with a medical fine-tuned model~(e.g., BioMedLM~\cite{bolton2022} or
MedPaLM) would improve report quality.
(iv)~The system currently supports episodic upload-and-analyze workflows only, not
continuous wearable monitoring streams.
(v)~The MIT-BIH heartbeat dataset exhibits strong class imbalance~(Normal: $82.8\%$);
although class-weighted loss mitigates this, minority-class performance~(Fusion, $80.8\%$~F1)
remains below the majority-class baseline.
(vi)~Heart disease predictions are based on UCI Cleveland clinical features and may
not generalize to patient populations from other geographic or ethnic cohorts without
re-calibration.

\subsection{Future Directions}

Planned extensions include:
(i)~HL7~FHIR API integration for EHR-connected deployment;
(ii)~federated learning across hospital systems to improve generalization without
centralizing patient data;
(iii)~temporal series analysis for continuous arrhythmia monitoring from wearable ECG
patches;
(iv)~domain-adapted clinical LLM substitution (MedPaLM or BioMedLM);
(v)~prospective clinical validation study with electrophysiology-confirmed labels;
(vi)~extension of the 1D-CNN heartbeat classifier to the full MIT-BIH raw signal
database~\cite{moody2001} with patient-adaptive fine-tuning; and
(vii)~integration of an external risk-factor API (e.g., ACC/AHA ASCVD Risk Calculator)
to augment the heart disease prediction pipeline with population-level risk scores.

%% ══════════════════════════════════════════════════════════════════════════════
%% X. CONCLUSION
%% ══════════════════════════════════════════════════════════════════════════════
\section{Conclusion}
\label{sec:conclusion}

This paper presented DeepCardio-RAG, an expanded multi-modal agentic AI framework
for cardiac arrhythmia detection, heart disease risk prediction, and clinical decision
support, designed for consumer health electronics deployment.
The six principal contributions address complementary diagnostic challenges:
CardioFusion multi-modal fusion~(AUC-ROC $0.941$),
arrhythmia video 3D-CNN~($89.7\%$ sensitivity),
LangChain+LangGraph RAG agent~($3.1\,\text{s}$ full-pipeline latency),
Tabular BERT+MoE arthritis predictor~($87.3\%$ F1),
UCI Cleveland heart disease predictor~($88.6\%$ accuracy, $0.934$ AUC-ROC), and
MIT-BIH ECG heartbeat classifier~($97.8\%$ weighted F1, $0.982$ AUC-ROC).
The two new Kaggle dataset integrations demonstrate that the unified Tabular BERT+MoE
and 1D-CNN architectures generalize across heterogeneous cardiovascular datasets
within a single extensible framework.
The open-source codebase, comprehensive REST API, and reproducible training
pipelines provide a platform for the research community to extend DeepCardio-RAG
to additional clinical domains.
DeepCardio-RAG demonstrates that agentic AI combining specialist neural models with
retrieval-augmented generation represents a productive architectural direction for
consumer health electronics, enabling transparent, multi-evidence, and traceable
clinical decision support.

%% ── Acknowledgment ────────────────────────────────────────────────────────────
\section*{Acknowledgment}

The author thanks the EchoNet-Dynamic team at Stanford AIMI for making the
echocardiogram video dataset publicly available, the PhysioNet community for
open-access cardiac signal datasets including MIT-BIH~\cite{moody2001},
the UCI Machine Learning Repository for the Cleveland Heart Disease
dataset~\cite{uci_heart,detrano1989}, and the Kaggle data contributors
(\textit{johnsmith88} and \textit{sadmansakib7}) whose curated distributions
enabled rapid prototyping.

%% ══════════════════════════════════════════════════════════════════════════════
%% REFERENCES  (BibTeX — also saved as references.bib)
%% ══════════════════════════════════════════════════════════════════════════════
\balance  % Balance last-page columns
\bibliographystyle{IEEEtran}
\bibliography{references}

%% ── If not using BibTeX, use thebibliography: ─────────────────────────────────
%\begin{thebibliography}{18}
%
%\bibitem{who2023}
%World Health Organization, ``Cardiovascular diseases (CVDs),'' WHO Fact Sheet, 2023.
%[Online]. Available: \url{https://www.who.int/news-room/fact-sheets/detail/cardiovascular-diseases-(cvds)}
%
%\bibitem{steinhubl2013}
%S.~R. Steinhubl, E.~D. Muse, and E.~J. Topol,
%``Can mobile health technologies transform health care?''
%\textit{JAMA}, vol.~310, no.~22, pp.~2395--2396, 2013.
%
%\bibitem{ramachandram2017}
%D.~Ramachandram and G.~W. Taylor,
%``Deep multimodal learning: A survey on recent advances and trends,''
%\textit{IEEE Signal Process.\ Mag.}, vol.~34, no.~6, pp.~96--108, 2017.
%
%\bibitem{lewis2020}
%P.~Lewis \textit{et al.},
%``Retrieval-augmented generation for knowledge-intensive NLP tasks,''
%in \textit{Proc.\ NeurIPS}, 2020, pp.~9459--9474.
%
%\bibitem{ouyang2020}
%D.~Ouyang \textit{et al.},
%``Video-based AI for beat-to-beat assessment of cardiac function,''
%\textit{Nature}, vol.~580, pp.~252--256, 2020.
%
%\bibitem{hannun2019}
%A.~Y. Hannun \textit{et al.},
%``Cardiologist-level arrhythmia detection and classification in ambulatory
% electrocardiograms using a deep neural network,''
%\textit{Nature Med.}, vol.~25, pp.~65--69, 2019.
%
%\bibitem{sun2024}
%C.~Sun, C.~Liu, X.~Wang, Y.~Liu, and S.~Zhao,
%``Coronary artery disease detection based on a novel multi-modal deep-coding method using ECG and PCG signals,''
%\textit{Sensors}, vol.~24, no.~21, art.~6939, 2024.
%
%\bibitem{gu2022}
%Y.~Gu \textit{et al.},
%``Domain-specific language model pretraining for biomedical NLP,''
%\textit{ACM Trans.\ Comput.\ Healthcare}, vol.~3, no.~1, pp.~1--23, 2022.
%
%\bibitem{zakka2024}
%C.~Zakka \textit{et al.},
%``Almanac: Retrieval-augmented language models for clinical medicine,''
%\textit{NEJM AI}, vol.~1, no.~2, 2024.
%
%\bibitem{langchain2022}
%Harrison Chase,
%``LangChain: Building Applications with LLMs through Composability,''
%GitHub, 2022. [Online]. Available: \url{https://github.com/langchain-ai/langchain}
%
%\bibitem{langgraph2024}
%LangChain AI,
%``LangGraph: Building Stateful, Multi-Actor Applications with LLMs,''
%GitHub, 2024. [Online]. Available: \url{https://github.com/langchain-ai/langgraph}
%
%\bibitem{mehandru2024}
%N.~Mehandru, B.~Y. Miao \textit{et al.},
%``Evaluating large language models as agents in the clinic,''
%\textit{npj Digital Med.}, vol.~7, art.~84, 2024.
%
%\bibitem{chen2016xgboost}
%T.~Chen and C.~Guestrin,
%``XGBoost: A scalable tree boosting system,''
%in \textit{Proc.\ KDD}, 2016, pp.~785--794.
%
%\bibitem{arik2021}
%S.~O. Arik and T.~Pfister,
%``TabNet: Attentive interpretable tabular learning,''
%in \textit{Proc.\ AAAI}, 2021, pp.~6679--6687.
%
%\bibitem{gorishniy2021}
%Y.~Gorishniy, I.~Rubachev \textit{et al.},
%``Revisiting deep learning models for tabular data,''
%in \textit{Proc.\ NeurIPS}, 2021, pp.~18932--18943.
%
%\bibitem{huang2020}
%Z.~Huang \textit{et al.},
%``TabTransformer: Tabular data modeling using contextual embeddings,''
%\textit{arXiv}:2012.06678, 2020.
%
%\bibitem{shazeer2017}
%N.~Shazeer \textit{et al.},
%``Outrageously large neural networks: The sparsely-gated mixture-of-experts layer,''
%in \textit{Proc.\ ICLR}, 2017.
%
%\bibitem{bolton2022}
%E.~Bolton \textit{et al.},
%``BioMedLM: A domain-specific large language model for biomedical text,''
%\textit{arXiv}:2301.09186, 2022.
%
%\bibitem{papineni2002}
%K.~Papineni \textit{et al.},
%``BLEU: A method for automatic evaluation of machine translation,''
%in \textit{Proc.\ ACL}, 2002, pp.~311--318.
%
%\bibitem{uci_heart}
%A.~Janosi \textit{et al.},
%``Heart Disease Dataset (UCI Cleveland + Hungarian + Swiss + VA),''
%Kaggle, contributed by johnsmith88, 2019.
%[Online]. Available: \url{https://www.kaggle.com/datasets/johnsmith88/heart-disease-dataset}
%
%\bibitem{detrano1989}
%R.~Detrano \textit{et al.},
%``International application of a new probability algorithm for the diagnosis of coronary artery disease,''
%\textit{Am.\ J.\ Cardiol.}, vol.~64, no.~5, pp.~304--310, 1989.
%
%\bibitem{moody2001}
%G.~B. Moody and R.~G. Mark,
%``The impact of the MIT-BIH Arrhythmia Database,''
%\textit{IEEE Eng.\ Med.\ Biol.\ Mag.}, vol.~20, no.~3, pp.~45--50, 2001.
%
%\bibitem{sadmansakib2021}
%S.~Sakib,
%``ECG Arrhythmia Classification Dataset (MIT-BIH derived, 187-feature segments),''
%Kaggle, 2021.
%[Online]. Available: \url{https://www.kaggle.com/datasets/sadmansakib7/ecg-arrhythmia-classification-dataset}
%
%\bibitem{de2004}
%P.~de~Chazal, M.~O'Dwyer, and R.~B. Reilly,
%``Automatic classification of heartbeats using ECG morphology and heartbeat interval features,''
%\textit{IEEE Trans.\ Biomed.\ Eng.}, vol.~51, no.~7, pp.~1196--1206, 2004.
%
%\bibitem{kiranyaz2015}
%S.~Kiranyaz, T.~Ince, and M.~Gabbouj,
%``Real-time patient-specific ECG classification by 1-D convolutional neural networks,''
%\textit{IEEE Trans.\ Biomed.\ Eng.}, vol.~63, no.~3, pp.~664--675, 2015.
%
%\bibitem{mohan2019}
%S.~Mohan, C.~Thirumalai, and G.~Srivastava,
%``Effective heart disease prediction using hybrid machine learning techniques,''
%\textit{IEEE Access}, vol.~7, pp.~81542--81554, 2019.
%
%\end{thebibliography}

\end{document}
